% Generated mechanically from the frozen brauer.tex target.
% Do not edit this generated file; edit the bound locale source instead.

% OK, start here.
%

\setcounter{chapter}{10}
\chapter{Brauer 群}



\phantomsection
\label{brauer-section-phantom}


\section{引言}
\label{brauer-section-introduction}

\noindent
一个参考资料是 Serre 在 Séminaire Cartan 上的讲义，见
\cite{Serre-Cartan}。Serre 又引用了
\cite{Deuring} 和 \cite{ANT}。我们改写了其中的一些证明，特别是使用
Rieffel 的一个巧妙论证来证明 Wedderburn 定理。
这一改写很可能并非改进；我们强烈建议读者阅读 Serre 的原始论述。


\section{非交换代数}
\label{brauer-section-algebras}

\noindent
设 $k$ 为域。本章所谓域 $k$ 上的一个{\it 代数} $A$，是指一个可能
非交换的环 $A$，连同一个环同态 $k \to A$，使得 $k$ 的像包含于
$A$ 的中心，且 $1$ 映到 $A$ 的单位元。所谓一个 {\it $A$-模}，
是指 $A$ 的单位元按恒等映射作用的右 $A$-模。

\begin{definition}
\label{brauer-definition-finite}
设 $A$ 为 $k$-代数。若 $\dim_k(A) < \infty$，则称 $A$ 为
{\it 有限维的}。此时记 $[A : k] = \dim_k(A)$。
\end{definition}

\begin{definition}
\label{brauer-definition-skew-field}
{\it 斜域}是一个带单位元 $1$ 的、可能非交换的环，其中
$1 \not = 0$，且每个非零元素都有乘法逆元。
\end{definition}

\noindent
一个斜域总是某个 $k$ 上的代数（例如可取它所含的素域作为 $k$）。
下文将使用如下事实：斜域上的任意模都是自由模，因为任一极大的
线性无关向量集构成一组基，而这样的集合由 Zorn 引理存在。

\begin{definition}
\label{brauer-definition-simple}
设 $A$ 为 $k$-代数。
若一个 $A$-模 $M$ 非零，且其仅有的 $A$-子模为 $0$ 和 $M$，
则称 $M$ 是{\it 单模}。
若 $A$ 的双边理想只有 $0$ 和 $A$，则称 $A$ 是{\it 单的}。
\end{definition}

\begin{definition}
\label{brauer-definition-central}
若 $A$ 的中心恰为 $k \to A$ 的像，则称 $k$-代数 $A$ 是
{\it 中心的}。
\end{definition}

\begin{definition}
\label{brauer-definition-opposite}
给定一个 $k$-代数 $A$，以 $A^{op}$ 表示将 $A$ 中乘法次序反转
所得的 $k$-代数。它称为{\it 反代数}。
\end{definition}




\section{Wedderburn 定理}
\label{brauer-section-wedderburn}

\noindent
以下巧妙论证可见于 Rieffel 的一篇论文，见
\cite{Rieffel}。这个证明不能再简单了（引自 Carl Faith 的评论）。

\begin{lemma}
\label{brauer-lemma-rieffel}
设 $A$ 是一个带 $1$ 的、可能非交换的环，且没有非平凡双边理想。
设 $M$ 是 $A$ 的一个非零右理想，并将 $M$ 视为右 $A$-模。
则 $A$ 与 $M$ 的双交换子相同。
\end{lemma}

\begin{proof}
令 $A' = \text{End}_A(M)$，于是 $M$ 是左 $A'$-模。
置 $A'' = \text{End}_{A'}(M)$（即 $M$ 的双交换子）。
我们将 $M$ 视为右 $A''$-模\footnote{这意味着，给定
$a'' \in A''$ 与 $m \in M$，有乘积 $m a'' \in M$。特别地，
$A''$ 中的乘法次序与把 $A''$ 的元素写成从左作用的自同态时
得到的次序相反。}。
令 $R : A \to A''$ 为满足 $mR(a) = ma$ 的自然同态。
由于 $R(1) = \text{id}_M$ 且 $A$ 没有非平凡双边理想，$R$ 是单射。
我们断言 $R(M)$ 是 $A''$ 的一个右理想。事实上，对
$a'' \in A''$ 和 $m$ 属于 $M$ 时有 $R(m)a'' = R(ma'')$，因为由
$M$ 的任一元素 $n$ 给出的左乘表示 $A'$ 的一个元素，从而对
所有属于 $M$ 的 $n$ 均有 $(nm)a'' = n(ma'')$。
最后，乘积理想 $AM$ 是双边理想，故 $A = AM$。于是
$R(A) = R(A)R(M)$，所以 $R(A)$ 是 $A''$ 的右理想。
但 $R(A)$ 含有 $A''$ 的单位元，故 $R(A) = A''$。
\end{proof}

\begin{lemma}
\label{brauer-lemma-simple-module}
设 $A$ 为 $k$-代数。若 $A$ 是有限维的，则
\begin{enumerate}
\item $A$ 有单模；
\item 任一非零模都含有单子模；
\item $A$ 上的单模在 $k$ 上是有限维的；并且
\item 若 $M$ 是单 $A$-模，则 $\text{End}_A(M)$ 是斜域。
\end{enumerate}
\end{lemma}

\begin{proof}
当然，由于 $A$ 自身是一个非零 $A$-模，(1) 可由 (2) 推出。
对于 (2)，任一作为 $k$-向量空间具有极小（有限）维数的子模都是
单模。这样的有限维子模存在，因为循环子模就是有限维的。
若 $M$ 是单模，则 $mA \subset M$ 是一个子模，由此得到 (3)。
$\text{End}_A(M)$ 的任一非零元素都是同构，故 (4) 成立。
\end{proof}

\begin{theorem}
\label{brauer-theorem-wedderburn}
\begin{slogan}
域上的有限维单代数是斜域上的矩阵代数。
\end{slogan}
设 $A$ 是一个有限维单 $k$-代数。则 $A$ 是某个有限维
$k$-代数 $K$ 上的矩阵代数，其中后者是斜域。
\end{theorem}

\begin{proof}
可以选择一个单子模 $M \subset A$，此时
$k$-代数 $K = \text{End}_A(M)$ 是斜域，见
引理 \ref{brauer-lemma-simple-module}。
由
引理 \ref{brauer-lemma-rieffel}
可知 $A = \text{End}_K(M)$。由于 $K$ 是斜域且
$M$ 有限生成（因为 $\dim_k(M) < \infty$），可知
$M$ 作为左 $K$-模是有限秩自由模。立即得到
$A \cong \text{Mat}(n \times n, K^{op})$。
\end{proof}






\section{代数的若干引理}
\label{brauer-section-lemmas}

\noindent
设 $A$ 为 $k$-代数，$B \subset A$ 为子代数。
{\it $B$ 在 $A$ 中的中心化子}是子代数
$$
C  = \{y \in A \mid xy = yx \text{ 对所有 }x \in B\}.
$$
它是一个 $k$-代数。

\begin{lemma}
\label{brauer-lemma-centralizer}
设 $A$、$A'$ 为 $k$-代数。设 $B \subset A$、$B' \subset A'$ 为
子代数，其中心化子分别为 $C$、$C'$。则
$B \otimes_k B'$ 在 $A \otimes_k A'$ 中的中心化子为
$C \otimes_k C'$。
\end{lemma}

\begin{proof}
以 $C'' \subset A \otimes_k A'$ 表示 $B \otimes_k B'$ 的
中心化子。显然 $C \otimes_k C' \subset C''$。反过来，
$C''$ 的每个元素都与 $B \otimes 1$ 交换，因而属于
$C \otimes_k A'$。类似地，$C'' \subset A \otimes_k C'$。故
$C'' \subset C \otimes_k A' \cap A \otimes_k C' = C \otimes_k C'$。
\end{proof}

\begin{lemma}
\label{brauer-lemma-center-csa}
设 $A$ 是有限维单 $k$-代数。则 $A$ 的中心 $k'$ 是 $k$ 的
有限域扩张。
\end{lemma}

\begin{proof}
由
定理 \ref{brauer-theorem-wedderburn}，
可写成 $A = \text{Mat}(n \times n, K)$，其中 $K$ 是
$k$ 上的有限维斜域。
由
引理 \ref{brauer-lemma-centralizer}，
$A$ 的中心为 $k \otimes_k k'$，其中 $k' \subset K$ 是 $K$ 的
中心。斜域的中心是域，结论得证。
\end{proof}

\begin{lemma}
\label{brauer-lemma-generate-two-sided-sub}
设 $V$ 为 $k$-向量空间。设 $K$ 是中心 $k$-代数且为斜域。
设 $W \subset V \otimes_k K$ 是一个双边 $K$-子向量空间。
则作为左 $K$-向量空间，$W$ 由 $W \cap (V \otimes 1)$ 生成。
\end{lemma}

\begin{proof}
令 $V' \subset V$ 为所有满足 $v \otimes 1 \in W$ 的
$v \in V$ 所生成的 $k$-子向量空间。于是
$V' \otimes_k K \subset W$，并且
$$
W/(V' \otimes_k K)  \subset  (V/V') \otimes_k K.
$$
若 $\overline{v} \in V/V'$ 是一个非零向量，且
$\overline{v} \otimes 1$ 属于 $W/(V' \otimes_k K)$，
则对 $\overline{v}$ 的一个提升 $v \in V$，有
$v \otimes 1 \in W$。这与 $V'$ 的构造矛盾。因此可将
$V$ 替换为 $V/V'$，将 $W$ 替换为 $W/(V' \otimes_k K)$；
只需证明当 $W$ 非零时，$W \cap (V \otimes 1)$ 非零。

\medskip\noindent
为此，取非零元素 $w \in W$，并将其写成
$w = \sum_{i = 1, \ldots, n} v_i \otimes k_i$，其中 $n$ 最小。
右乘 $k_1^{-1}$ 后，可设 $k_1 = 1$。
若 $n = 1$，则因 $v_1 \otimes 1 \in W$，结论成立。
若 $n > 1$，则对任意 $c \in K$，
$$
c w - w c = \sum\nolimits_{i = 2, \ldots, n} v_i \otimes (c k_i - k_i c) \in W
$$
而由 $n$ 的极小性，必有 $c k_i - k_i c = 0$。
这说明 $k_i$ 位于 $K$ 的中心；按假设该中心就是 $k$。
于是 $w = (v_1 + \sum k_i v_i) \otimes 1$，与 $n$ 的极小性矛盾。
\end{proof}

\begin{lemma}
\label{brauer-lemma-generate-two-sided-ideal}
设 $A$ 为 $k$-代数。设 $K$ 是中心 $k$-代数且为斜域。
则任一双边理想 $I \subset A \otimes_k K$ 都可写成
$J \otimes_k K$，其中 $J \subset A$ 是双边理想。
特别地，若 $A$ 是单的，则 $A \otimes_k K$ 也是单的。
\end{lemma}

\begin{proof}
令 $J = \{a \in A \mid a \otimes 1 \in I\}$。这是 $A$ 的一个
双边理想。由
引理 \ref{brauer-lemma-generate-two-sided-sub}
有 $I = J \otimes_k K$。
\end{proof}

\begin{lemma}
\label{brauer-lemma-matrix-algebras}
设 $R$ 是一个可能非交换的环，$n \geq 1$ 是整数。
令 $R_n = \text{Mat}(n \times n, R)$。
\begin{enumerate}
\item 函子 $M \mapsto M^{\oplus n}$ 与
$N \mapsto Ne_{11}$ 给出范畴之间的一对拟逆等价
$\text{Mod}_R \leftrightarrow \text{Mod}_{R_n}$。
\item $R_n$ 的任一双边理想都形如 $IR_n$，其中
$I$ 是 $R$ 的双边理想。
\item $R_n$ 的中心等于 $R$ 的中心。
\end{enumerate}
\end{lemma}

\begin{proof}
第 (1) 项是显然的。若 $J \subset R_n$ 是双边理想，则
$J = \bigoplus e_{ii}Je_{jj}$，且所有分量 $e_{ii}Je_{jj}$
彼此相等，并构成 $R$ 的一个双边理想 $I$。这证明了 (2)。
第 (3) 项显然。
\end{proof}

\begin{lemma}
\label{brauer-lemma-simple-module-unique}
设 $A$ 是有限维单 $k$-代数。
\begin{enumerate}
\item 在同构意义下，恰有一个单 $A$-模 $M$。
\item 任一有限维 $A$-模都是若干个单模副本的直和。
\item 两个有限维 $A$-模同构，当且仅当它们在 $k$ 上的维数相同。
\item 若 $A = \text{Mat}(n \times n, K)$，其中 $K$ 是 $k$ 的一个
有限维斜域扩张，则 $M = K^{\oplus n}$ 是单 $A$-模，且
$\text{End}_A(M) = K^{op}$。
\item 若 $M$ 是单 $A$-模，则 $L = \text{End}_A(M)$ 是 $k$ 上的
有限维斜域，并从左作用于 $M$；此外
$A = \text{End}_L(M)$，且 $A$ 与 $L$ 的中心相同。
并且 $[A : k] [L : k] = \dim_k(M)^2$。
\item 对于有限维 $A$-模 $N$，代数 $B = \text{End}_A(N)$ 是
(5) 中斜域 $L$ 上的矩阵代数。此外 $\text{End}_B(N) = A$。
\end{enumerate}
\end{lemma}

\begin{proof}
由
定理 \ref{brauer-theorem-wedderburn}，
可写成 $A = \text{Mat}(n \times n, K)$，其中 $K$ 是
$k$ 的一个有限维斜域扩张。由
引理 \ref{brauer-lemma-matrix-algebras}，
$A$-模范畴等价于 $K$-模范畴。由于 $K$ 上的每个模都是自由模，
(1)、(2)、(3) 成立。第 (4) 项成立，因为该等价将 $K$-模 $K$
变换为 $M = K^{\oplus n}$。在 (5) 中取 $M = K^{\oplus n}$，
可见 $L = K^{op}$。关于 $L = K^{op}$ 的中心的断言来自
引理 \ref{brauer-lemma-matrix-algebras}。
关于 $\text{End}_L(M)$ 的断言由 $M$ 的显式形式得到。
维数公式显然。第 (6) 项成立，因为 $N$ 同构于若干个单模副本的直和。
\end{proof}

\begin{lemma}
\label{brauer-lemma-tensor-simple}
设 $A$、$A'$ 是两个单 $k$-代数，其中一个在 $k$ 上有限维且中心。
则 $A \otimes_k A'$ 是单的。
\end{lemma}

\begin{proof}
设 $A'$ 在 $k$ 上有限维且中心。
由
定理 \ref{brauer-theorem-wedderburn}，
可写 $A' = \text{Mat}(n \times n, K')$。
此时 $K'$ 的中心为 $k$，由
引理 \ref{brauer-lemma-generate-two-sided-ideal}
可知 $A \otimes_k K'$ 是单的。
因此，由引理 \ref{brauer-lemma-matrix-algebras}，
$A \otimes_k A' = \text{Mat}(n \times n, A \otimes_k K')$ 是单的。
\end{proof}

\begin{lemma}
\label{brauer-lemma-tensor-central-simple}
域 $k$ 上两个有限维中心单代数的张量积仍是有限维、中心且单的。
\end{lemma}

\begin{proof}
合并引理 \ref{brauer-lemma-centralizer} 与 \ref{brauer-lemma-tensor-simple}。
\end{proof}

\begin{lemma}
\label{brauer-lemma-base-change}
设 $A$ 是 $k$ 上的有限维中心单代数，$k'/k$ 是域扩张。
则 $A' = A \otimes_k k'$ 是 $k'$ 上的有限维中心单代数。
\end{lemma}

\begin{proof}
合并引理 \ref{brauer-lemma-centralizer} 与 \ref{brauer-lemma-tensor-simple}。
\end{proof}

\begin{lemma}
\label{brauer-lemma-inverse}
设 $A$ 是 $k$ 上的有限维中心单代数。则
$A \otimes_k A^{op} \cong \text{Mat}(n \times n, k)$，
其中 $n = [A : k]$。
\end{lemma}

\begin{proof}
由引理 \ref{brauer-lemma-tensor-central-simple}，
代数 $A \otimes_k A^{op}$ 是单的。因此映射
$$
A \otimes_k A^{op} \longrightarrow \text{End}_k(A),\quad
a \otimes a' \longmapsto (x \mapsto axa')
$$
是单射。箭头两端维数相同，故结论成立。
\end{proof}





\section{域的 Brauer 群}
\label{brauer-section-brauer}

\noindent
设 $k$ 为域。考虑 $k$ 上两个有限维中心单代数 $A$ 与 $B$。
若存在 $n, m > 0$，使得作为 $k$-代数有
$\text{Mat}(n \times n, A) \cong \text{Mat}(m \times m, B)$，
则称 $A$ 与 $B$ {\it 相似}。

\begin{lemma}
\label{brauer-lemma-similar}
关于相似性，有：
\begin{enumerate}
\item 相似性在 $k$ 上有限维中心单代数的同构类集合上定义了一个
等价关系。
\item 每个相似类都含有唯一的（在同构意义下）$k$ 上有限维中心斜域。
\item 若 $A = \text{Mat}(n \times n, K)$ 且
$B = \text{Mat}(m \times m, K')$，其中 $K$、$K'$ 是 $k$ 上的
有限维中心斜域，则 $A$ 与 $B$ 相似，当且仅当
$K \cong K'$（作为 $k$-代数）。
\end{enumerate}
\end{lemma}

\begin{proof}
注意，由 Wedderburn 定理（定理 \ref{brauer-theorem-wedderburn}），
任一有限维中心单代数都可写成某个有限维中心斜域上的矩阵代数。
因此只需证明第三个断言。为此，只需证明若
$A = \text{Mat}(n \times n, K) \cong \text{Mat}(m \times m, K') = B$，
则 $K \cong K'$。事实上，对 $A$ 的一个单模 $M$，由
引理 \ref{brauer-lemma-simple-module-unique}
有 $\text{End}_A(M) = K^{op}$。
所以 $A \cong B$ 蕴含 $K^{op} \cong (K')^{op}$，结论得证。
\end{proof}

\noindent
给定两个有限维中心单 $k$-代数 $A$、$B$，其张量积
$A \otimes_k B$ 仍是这样的代数，见
引理 \ref{brauer-lemma-tensor-central-simple}。
此外，若 $A$ 与 $A'$ 相似，则 $A \otimes_k B$ 与
$A' \otimes_k B$ 相似，因为张量积与取矩阵代数可交换。
因此张量积在有限维中心单代数的等价类上定义了一个显然满足
结合律和交换律的运算。最后，
引理 \ref{brauer-lemma-inverse}
表明 $A \otimes_k A^{op}$ 同构于矩阵代数，也就是说，
$A \otimes_k A^{op}$ 属于 $k$ 的相似类。
由此得到一个阿贝尔群。

\begin{definition}
\label{brauer-definition-brauer-group}
设 $k$ 为域。$k$ 的 {\it Brauer 群}是上述有限维中心单
$k$-代数的相似类所成的阿贝尔群，记作 $\text{Br}(k)$。
\end{definition}

\noindent
对任意域同态 $k \to k'$，由引理 \ref{brauer-lemma-base-change}
得到群同态
$$
\text{Br}(k) \longrightarrow \text{Br}(k'),\quad
A \longmapsto A \otimes_k k'
$$
。换言之，$\text{Br}(-)$ 是从域范畴到阿贝尔群范畴的函子。
注意，一个域的 Brauer 群为零，当且仅当每个有限维中心斜域扩张
$k \subset K$ 都是平凡的。

\begin{lemma}
\label{brauer-lemma-brauer-algebraically-closed}
代数闭域的 Brauer 群为零。
\end{lemma}

\begin{proof}
设 $k \subset K$ 为有限维中心斜域扩张。
对任意元素 $x \in K$，子环 $k[x] \subset K$ 是一个交换的、
有限维且整的 $k$-子代数，因而是域，见
《交换代数》，引理 \ref{algebra-lemma-integral-over-field}。
由于 $k$ 代数闭，得到 $k[x] = k$。因 $x$ 任意，遂有 $k = K$。
\end{proof}

\begin{lemma}
\label{brauer-lemma-dimension-square}
设 $A$ 是域 $k$ 上的有限维中心单代数，则 $[A : k]$ 是完全平方数。
\end{lemma}

\begin{proof}
这是因为由
引理 \ref{brauer-lemma-brauer-algebraically-closed}，
$A \otimes_k \overline{k}$ 是 $\overline{k}$ 上的矩阵代数。
\end{proof}




\section{Skolem--Noether 定理}
\label{brauer-section-skolem-noether}



\begin{theorem}
\label{brauer-theorem-skolem-noether}
设 $A$ 是有限维中心单 $k$-代数，$B$ 是单 $k$-代数，并设
$f, g : B \to A$ 是两个 $k$-代数同态。则存在可逆元素
$x \in A$，使得对所有 $b \in B$ 都有
$f(b) = xg(b)x^{-1}$。
\end{theorem}

\begin{proof}
选择一个单 $A$-模 $M$，并令 $L = \text{End}_A(M)$。
由引理 \ref{brauer-lemma-simple-module} 和
\ref{brauer-lemma-simple-module-unique}，
$L$ 是中心为 $k$ 的斜域，并从左作用于 $M$。
此时 $M$ 上有两个 $B \otimes_k L^{op}$-模结构，分别定义为
$m \cdot_1 (b \otimes l) = lmf(b)$ 和
$m \cdot_2 (b \otimes l) = lmg(b)$。
由引理 \ref{brauer-lemma-tensor-simple}，
$k$-代数 $B \otimes_k L^{op}$ 是单的。由于 $B$ 是单的，
$k$-代数同态 $B \to A$ 的存在性蕴含 $B$ 有限维。
因此 $B \otimes_k L^{op}$ 是有限维单代数；由
引理 \ref{brauer-lemma-simple-module-unique}，
$M$ 上这两个 $B \otimes_k L^{op}$-模结构同构。
所以存在交织这两个作用的 $\varphi : M \to M$。
特别地，$\varphi$ 属于 $L$ 的交换子，这说明 $\varphi$
是由某个 $x \in A$ 给出的乘法，见
引理 \ref{brauer-lemma-simple-module-unique}。展开定义可见，$x$
正是所求解。
\end{proof}

\begin{lemma}
\label{brauer-lemma-automorphism-inner}
设 $A$ 是有限维中心单 $k$-代数。$A$ 的任一自同构都是内自同构。
特别地，$\text{Mat}(n \times n, k)$ 的任一自同构都是内自同构。
\end{lemma}

\begin{proof}
应用 Skolem--Noether 定理（定理 \ref{brauer-theorem-skolem-noether}）。
\end{proof}



\section{中心化子定理}
\label{brauer-section-centralizer}



\begin{theorem}
\label{brauer-theorem-centralizer}
设 $A$ 是 $k$ 上的有限维中心单代数，$B$ 是 $A$ 的单子代数。
则
\begin{enumerate}
\item $B$ 在 $A$ 中的中心化子 $C$ 是单的；
\item $[A : k] = [B : k][C : k]$；并且
\item $C$ 在 $A$ 中的中心化子是 $B$。
\end{enumerate}
\end{theorem}

\begin{proof}
在本证明中，我们将自由使用
引理 \ref{brauer-lemma-simple-module-unique} 的结论。
选择一个单 $A$-模 $M$，并令 $L = \text{End}_A(M)$。
则 $L$ 是中心为 $k$ 的斜域，从左作用于 $M$，并且
$A = \text{End}_L(M)$。
于是 $M$ 是右 $B \otimes_k L^{op}$-模，且
$C = \text{End}_{B \otimes_k L^{op}}(M)$。
由引理 \ref{brauer-lemma-tensor-simple}，代数
$B \otimes_k L^{op}$ 是单的，故 $C$ 是单的（再次应用
引理 \ref{brauer-lemma-simple-module-unique}）。

\medskip\noindent
写成 $B \otimes_k L^{op} = \text{Mat}(m \times m, K)$，
其中 $K$ 是 $k$ 上的有限维斜域。若 $M$ 同构于单
$B \otimes_k L^{op}$-模 $K^{\oplus m}$ 的 $n$ 个副本的直和，
则 $C = \text{Mat}(n \times n, K^{op})$（再次应用该引理）。
因此 $\dim_k(M) = nm [K : k]$，
$[B : k] [L : k] = m^2 [K : k]$，
$[C : k] = n^2 [K : k]$，且
$[A : k] [L : k] = \dim_k(M)^2$（仍由该引理）。
由此得到 (2)。

\medskip\noindent
对 $C \subset A$ 应用 (2)，可知 (3) 成立：若 $C'$ 是
$C$ 在 $A$ 中的中心化子，则 $[B : k] = [C' : k]$，
而显然 $B \subset C')$。
\end{proof}

\begin{lemma}
\label{brauer-lemma-when-tensor-is-equal}
设 $A$ 是 $k$ 上的有限维中心单代数，$B$ 是 $A$ 的单子代数。
若 $B$ 是中心 $k$-代数，则 $A = B \otimes_k C$，其中
$C$ 是 $B$ 在 $A$ 中的（中心单）中心化子。
\end{lemma}

\begin{proof}
由定理 \ref{brauer-theorem-centralizer}，
$\dim_k(A) = \dim_k(B \otimes_k C)$。
由引理 \ref{brauer-lemma-tensor-simple}，该张量积是单的。
因此自然映射 $B \otimes_k C \to A$ 是单射，从而是同构。
\end{proof}

\begin{lemma}
\label{brauer-lemma-self-centralizing-subfield}
设 $A$ 是 $k$ 上的有限维中心单代数。若 $K \subset A$ 是子域，
则以下条件等价：
\begin{enumerate}
\item $[A : k] = [K : k]^2$；
\item $K$ 等于它自身的中心化子；并且
\item $K$ 是极大交换子环。
\end{enumerate}
\end{lemma}

\begin{proof}
定理 \ref{brauer-theorem-centralizer} 表明 (1) 与 (2) 等价。
显然，(3) 与 (2) 等价。
\end{proof}

\begin{lemma}
\label{brauer-lemma-maximal-subfield}
\begin{slogan}
有限维中心斜域的维数等于其任一极大子域维数的平方。
\end{slogan}
设 $A$ 是 $k$ 上的有限维中心斜域。
则每个极大子域 $K \subset A$ 都满足
$[A : k] = [K : k]^2$。
\end{lemma}

\begin{proof}
这是引理 \ref{brauer-lemma-self-centralizing-subfield} 的特殊情形。
\end{proof}





\section{分裂域}
\label{brauer-section-splitting}



\begin{definition}
\label{brauer-definition-splitting}
设 $A$ 是有限维中心单 $k$-代数。若 $A \otimes_k k'$ 是
$k'$ 上的矩阵代数，则称域扩张 $k'/k$ {\it 分裂} $A$，
或称 $k'$ 是 $A$ 的一个{\it 分裂域}。
\end{definition}

\noindent
等价地说，$A$ 的类在映射
$\text{Br}(k) \to \text{Br}(k')$ 下映为零。

\begin{theorem}
\label{brauer-theorem-splitting}
设 $A$ 是有限维中心单 $k$-代数，$k'/k$ 是有限域扩张。
以下条件等价：
\begin{enumerate}
\item $k'$ 分裂 $A$；并且
\item 存在一个与 $A$ 相似的有限维中心单代数 $B$，使得
$k' \subset B$ 且 $[B : k] = [k' : k]^2$。
\end{enumerate}
\end{theorem}

\begin{proof}
假设 (2) 成立。只需证明 $B \otimes_k k'$ 是矩阵代数。
已知 $B \otimes_k B^{op} \cong \text{End}_k(B)$。
由引理 \ref{brauer-lemma-self-centralizing-subfield}，$k'$ 是
$B^{op}$ 中 $k'$ 的中心化子，故 $B \otimes_k k'$ 是
$k \otimes k'$ 在
$B \otimes_k B^{op} = \text{End}_k(B)$ 中的中心化子。
当然，这个中心化子就是 $\text{End}_{k'}(B)$，其中借助嵌入
$k' \to B$ 将 $B$ 视为 $k'$-向量空间。结论随即成立。

\medskip\noindent
假设 (1) 成立。这意味着对某个 $k'$-向量空间 $V$，存在同构
$A \otimes_k k' \cong \text{End}_{k'}(V)$。
令 $B$ 为 $A$ 在 $\text{End}_k(V)$ 中的交换子。注意
$k'$ 包含于 $B$。由
引理 \ref{brauer-lemma-when-tensor-is-equal}，
$A$ 与 $B$ 的类在 $\text{Br}(k)$ 中之和为零。
由
定理 \ref{brauer-theorem-centralizer}
中的维数公式，
$$
[B : k] [A : k] =
\dim_k(V)^2 =
[k' : k]^2 \dim_{k'}(V)^2 =
[k' : k]^2 [A : k].
$$
因此 $[B : k] = [k' : k]^2$。这样，我们已对 $A$ 的 Brauer
类的相反类证明了结论。但 $k'$ 分裂 $A$ 的 Brauer 类，当且仅当
它分裂反代数的 Brauer 类，故结论仍然成立。
\end{proof}

\begin{lemma}
\label{brauer-lemma-maximal-subfield-splits}
$k$ 上有限维中心斜域 $K$ 的任一极大子域都是 $K$ 的分裂域。
\end{lemma}

\begin{proof}
合并引理 \ref{brauer-lemma-maximal-subfield} 与
定理 \ref{brauer-theorem-splitting}。
\end{proof}

\begin{lemma}
\label{brauer-lemma-splitting-field-degree}
设 $K$ 是 $k$ 上的有限维中心斜域，并令 $d^2 = [K : k]$。
对 $K$ 的任一有限分裂域 $k'$，次数 $[k' : k]$ 可被 $d$ 整除。
\end{lemma}

\begin{proof}
由定理 \ref{brauer-theorem-splitting}，存在 Brauer 类与 $K$ 相同的
有限维中心单代数 $B$，使得 $[B : k] = [k' : k]^2$。
由
引理 \ref{brauer-lemma-similar}，
对某个 $n$ 有 $B = \text{Mat}(n \times n, K)$。
于是 $[k' : k]^2 = n^2d^2$，结论得证。
\end{proof}

\begin{proposition}
\label{brauer-proposition-separable-splitting-field}
设 $K$ 是 $k$ 上的有限维中心斜域。存在极大子域
$k \subset k' \subset K$，使得它在 $k$ 上可分。
% P02-E0001: frozen English source says “spitting field” at brauer.tex:711.
特别地，每个 Brauer 类都有一个有限可分分裂域。
\end{proposition}

\begin{proof}
由于每个 Brauer 类都可由 $k$ 上的有限维中心斜域表示，
第二个断言由第一个断言和
引理 \ref{brauer-lemma-maximal-subfield-splits}
立即得到。

\medskip\noindent
为证明第一个断言，设已给定可分子域 $k' \subset K$。
则 $k'$ 在 $K$ 中的中心化子 $K'$ 以 $k'$ 为中心，问题归结为
寻找 $K'$ 的一个在 $k'$ 上可分的极大子域。因此只需证明：
若 $k \not = K$，则存在 $x \in K$、$x \not \in k$，且该元素
在 $k$ 上可分。在特征零时，这是显然的。因此可设 $k$ 的特征为
$p > 0$。若基域 $k$ 有限，结论同样显然（因为有限域的扩张总是
可分的）。故可设 $k$ 是正特征无限域。

\medskip\noindent
% P02-E0002: source statement quantifies over all of K although elements of k are separable over k.
反设 $K$ 中没有元素在 $k$ 上可分，以求矛盾。由
《域》，第 \ref{fields-section-algebraic} 节
中的讨论，这意味着任意 $x \in K$ 的最小多项式都形如
$T^q - a$，其中 $q$ 是 $p$ 的幂且 $a \in k$。
显然，$K$ 中每个元素的最小多项式，其次数都满足
$\leq \dim_k(K)$，故存在一个固定的 $p$ 的幂 $q$，使得对所有
$x \in K$ 都有 $x^q \in k$。

\medskip\noindent
考虑如下映射：
$$
(-)^q : K \longrightarrow K
$$
取 $K$ 的一组 $k$-基 $\{a_1, \ldots, a_n\}$，其中
$a_1 = 1$，用这组基将该映射写成
$$
(\sum x_i a_i)^q = \sum f_i(x_1, \ldots, x_n)a_i.
$$
由于 $K$ 上的乘法是 $k$-双线性的，每个 $f_i$ 都是
$x_1, \ldots, x_n$ 的多项式（细节从略）。
上述 $q$ 的选择以及 $k$ 为无限域这一事实表明，对 $i \geq 2$，
$f_i$ 恒为零。因此，将 $k$ 扩张到其代数闭包
$\overline{k}$ 后，它仍恒为零。但代数
$K \otimes_k \overline{k}$ 是矩阵代数（例如由
引理 \ref{brauer-lemma-base-change} 与
\ref{brauer-lemma-brauer-algebraically-closed}），
这意味着存在某些元素，其 $q$ 次幂不是中心元素
（例如 $e_{11}$）。这就是所需的矛盾。
\end{proof}

\noindent
以上结果使我们能够如下刻画有限维中心单代数。

\begin{lemma}
\label{brauer-lemma-finite-central-simple-algebra}
设 $k$ 为域。对 $k$-代数 $A$，以下条件等价：
\begin{enumerate}
\item $A$ 是有限维中心单 $k$-代数；
\item $A$ 是有限维 $k$-向量空间，$k$ 是 $A$ 的中心，且
$A$ 没有非平凡双边理想；
\item 存在 $d \geq 1$，使得
$A \otimes_k \bar k \cong \text{Mat}(d \times d, \bar k)$；
\item 存在 $d \geq 1$，使得
$A \otimes_k k^{sep} \cong \text{Mat}(d \times d, k^{sep})$；
\item 存在 $d \geq 1$ 及有限 Galois 扩张 $k'/k$，使得
$A \otimes_k k' \cong \text{Mat}(d \times d, k')$；
\item 存在 $n \geq 1$ 及 $k$ 上的有限维中心斜域 $K$，使得
$A \cong \text{Mat}(n \times n, K)$。
\end{enumerate}
整数 $d$ 称为 $A$ 的{\it 次数}。
\end{lemma}

\begin{proof}
(1) 与 (2) 的等价性直接来自定义，见
Section \ref{brauer-section-algebras}。
设 (1) 成立。由
命题 \ref{brauer-proposition-separable-splitting-field}，
$A$ 有一个可分分裂域 $k \subset k'$。
当然，$k'/k$ 的一个 Galois 闭包也是分裂域。
所以 (1) 蕴含 (5)。显然，(5) $\Rightarrow$ (4)
$\Rightarrow$ (3)。设 (3) 成立。则
$A \otimes_k \overline{k}$ 是有限维中心单
$\overline{k}$-代数，例如可由
引理 \ref{brauer-lemma-matrix-algebras} 得到。
这显然蕴含 $A$ 是有限维中心单 $k$-代数。
最后，(1) 与 (6) 的等价性就是 Wedderburn 定理，见
定理 \ref{brauer-theorem-wedderburn}。
\end{proof}
